\documentclass[10pt,a4paper]{amsart}
\usepackage[margin=19mm]{geometry}
\usepackage{amsmath,amssymb,mathtools}
\usepackage[scr=rsfs,cal=euler]{mathalfa}
\usepackage{xfrac}
\usepackage[unicode,hidelinks,bookmarks=false]{hyperref}
\newcommand{\ie}{i.e.\ }
\newcommand{\cf}{cf.\ }
\newcommand{\resp}{respectively\ }
\newcommand{\Subsection}{\subsection}
\providecommand{\nobreakdash}{\nobreak\hbox{-}}
\setlength{\emergencystretch}{2em}
\multlinegap0pt
\usepackage{bm}
\usepackage{bbm}
\usepackage{dsfont}
\usepackage{xspace}
\usepackage{cancel}
\usepackage{mathtools}
\usepackage{cleveref}
\usepackage{amsthm}
\makeatletter
\@ifundefined{theorem}{\newtheorem{theorem}{Theorem}}{}
\@ifundefined{lemma}{\newtheorem{lemma}[theorem]{Lemma}}{}
\makeatother
\newcommand{\V}{\mathcal{V}}
\newcommand{\pa}{\text{par}}
\title[Greedy Multi-Path Block Verification for Faster Decoding in ]{Greedy Multi-Path Block Verification for Faster Decoding in Speculative Sampling}
\author{Rahul Thomas, Arka Pal}
\date{Source: arXiv:2602.16961v1 (CC BY 4.0)}
\begin{document}
\maketitle

\noindent\emph{Source and licence.} Greedy Multi-Path Block Verification for Faster Decoding in Speculative Sampling. Version arXiv:2602.16961v1 (CC BY 4.0), distributed under the Creative Commons Attribution 4.0 International licence, as recorded in the arXiv licence metadata for this exact version and shown on the abstract page https://arxiv.org/abs/2602.16961v1. Full rights notice: licenses/arxiv-2602.16961-CC-BY-4.0.txt.\par\smallskip

\noindent\textit{Editorial scope.} Counted unit: the Lemma whose source label is \texttt{None} in the article (source0.tex lines 705--713, source label \texttt{None}), delivered together with the article's own complete original proof of it (source0.tex lines 715--731, one \texttt{proof} environment of 17 lines). No mathematics is written, repaired or re-derived: statement and proof are the article's own text line for line. Required context retained uncounted: lem:prefix (lines 654--663); lem:target (lines 679--685). Every definition, display and label of those lines is retained unchanged. Omitted from this package and not redistributed: 1--653, 664--678, 686--704, 714--714, 732--1149. Nothing inside a retained excerpt is omitted or altered: every retained body is delivered exactly as the article's lines are written, so no omission receipt of any kind accompanies this package. Citations: the retained excerpts carry no citation command at all, so no bibliography entry of the article is transcribed and no reference is redistributed. Reference adaptation: none was needed and none was made; every reference inside the delivered unit resolves inside it, so no unresolved reference or citation is produced. Printed slips kept literal: no printed word, symbol, hypothesis or number of the counted statement or of its proof was altered. Layout and class: the article's own class is \texttt{article} with packages including microtype, graphicx, subcaption, booktabs, hyperref, icml2026, amsmath, amssymb, mathtools, amsthm, cleveref, multirow, todonotes; the wrapper is the lane's pinned \texttt{amsart} editorial wrapper at 10pt on A4, which restates the theorem-like environments the retained text needs and adds the article's own macro definitions that the retained text needs (\texttt{\textbackslash V}, \texttt{\textbackslash pa}), with extra wrapper packages (bm, bbm, dsfont, xspace, cancel, mathtools, cleveref). The delivered numbering is the unit's own. Assets: no retained span contains a figure environment, a graphics-inclusion command, a PDF-page inclusion, a table or a TikZ picture; no image file is copied into this package, the package declares no asset dependency and no image file is redistributed with it. Licence: version arXiv:2602.16961v1 of this article is distributed under the Creative Commons Attribution 4.0 International licence, as recorded in the arXiv licence metadata for this exact version and shown on the abstract page https://arxiv.org/abs/2602.16961v1. A full rights notice is delivered at licenses/arxiv-2602.16961-CC-BY-4.0.txt. Characters: every retained excerpt is pure ASCII, so the delivered engine, XeLaTeX with its default Latin Modern text font, reports no missing character.
\begin{lemma}[Prefix-Matching]
    Construct the bipartite graph $G$ with left vertices $V^{\leq L+1} \setminus \V^0$, right vertices $(\V^L)^K$, and edges $E$ consisting of all pairs \begin{equation}
        \left(a_{1:i},\left(a^{(1)}_{1:L},\ldots,a^{(K)}_{1:L}\right)\right) \in \left(\V^{\leq L+1}\setminus \V^0\right) \times \left(\V^L\right)^K, \quad a_{1:i-1} \in \left\{a^{(1)}_{1:i-1},\ldots,a^{(K)}_{1:i-1}\right\}.
    \end{equation} The values $R_\Phi(a_{1:i})$ for $a_{1:i} \in \V^{\leq L+1}\setminus \V^0$ are feasible for a $K$-path verification algorithm $\Phi$ that satisfies prefix-matching if and only if they are nonnegative, sum to one, and 
    \begin{equation}
        \sum_{a_{1:i} \in S} R_\Phi(a_{1:i}) \leq \sum_{\left(a^{(1)}_{1:L},\ldots,a^{(K)}_{1:L}\right) \in N(S)}  \prod_{k=1}^K q\left(a^{(k)}_{1:L}\right) \quad \forall S \subseteq V^{\leq L+1} \setminus \V^0,
    \end{equation}
    where $N(\cdot)$ denotes a neighborhood in $G$.
    \label{lem:prefix}
\end{lemma}

\begin{lemma}[Target-Matching]
    The values $R_\Phi(a_{1:i})$ for $a_{1:i} \in \V^{\leq L+1}\setminus \V^0$ are feasible for a $K$-path verification algorithm $\Phi$ that satisfies prefix-matching if and only if they are nonnegative, sum to one, and
    \begin{equation}
        \sum_{i=0}^{L} \frac{R_\Phi(a_{1:i+1})}{p(a_{1:i+1})} = 1\quad \forall a_{1:L+1} \in \V^{L+1}.
    \end{equation}
    \label{lem:target}
\end{lemma}

\begin{lemma}[Unsimplified Multi-Path LP]
    The values $R_\Phi(a_{1:i})$ for $a_{1:i} \in \V^{\leq L+1}\setminus \V^0$ are feasible for a valid $K$-path verification algorithm $\Phi$ if and only if they are nonnegative, sum to one, and are feasible in
    \begin{align}
        \max & \sum_{i=1}^{L+1} \sum_{a_{1:i} \in \V^{i}} i R_\Phi(a_{1:i}),\\
        \text{s.t.}\quad & \sum_{i=0}^{L} \frac{R_\Phi(a_{1:i+1})}{p(a_{1:i+1})} = 1\quad \forall a_{1:L+1} \in \V^{L+1},\\
        & \sum_{a_{1:i} \in T} \sum_{j=i+1}^{L+1} \sum_{a_{i+1:j} \in \V^{j-i}} R_\Phi(a_{1:j})  \leq 1 - \left(1 - \sum_{a_{1:i} \in T} q(a_{1:i})\right)^K \quad \forall T \in \mathcal{A}(\V^{\leq L}).
    \end{align}
    Furthermore, the objective value is precisely the block efficiency $\mathbb{E}[\tau+1]$ for $\Phi$.
\end{lemma}

\begin{proof}
First, fix a subset $S \subseteq \V^{\leq L+1} \setminus \V^0$. The neighborhood $N(S) \subseteq (\V^L)^K$ induced by the graph $G$ in \cref{lem:prefix} is the set of $K$-tuples of length-$L$ paths where at least one path has a prefix in $\pa(S)$. So, the complement of $N(S)$ consists of all $K$-tuples of length-$L$ paths where no path has a prefix in $\pa(S)$. Now, note that a path has a prefix in $\pa(S)$ if and only if it has a prefix in $\min(\pa(S))$. Because $\min(\pa(S))$ is an antichain (i.e. no two distinct elements in it are a prefix of the same path), disjointness shows that the probability that a length $L$-path sampled from $q(\cdot)$ has a prefix in $\min(\pa(S))$ is $\sum_{a_{1:i} \in \min(\pa(S))} q(a_{1:i})$. Using these observations, we simplify:
\begin{align}
    \sum_{\left(a^{(1)}_{1:L},\ldots,a^{(K)}_{1:L}\right) \in N(S)}  \prod_{k=1}^K q\left(a^{(k)}_{1:L}\right)
    &= 1 - \sum_{\left(a^{(1)}_{1:L},\ldots,a^{(K)}_{1:L}\right) \not \in N(S)}\prod_{k=1}^K q\left(a^{(k)}_{1:L}\right) \\
    &= 1 - \left(1 - \sum_{a_{1:i} \in \min(\pa(S))} q(a_{1:i})\right)^K.
\end{align}
Thus, the prefix-matching inequalities from \cref{lem:prefix} become:
\begin{equation}
    \sum_{a_{1:i} \in S} R_\Phi(a_{1:i}) \leq 1 - \left(1 - \sum_{a_{1:i} \in \min(\pa(S))} q(a_{1:i})\right)^K \quad \forall S \subseteq V^{\leq L+1} \setminus \V^0.
\end{equation}
We now make the key observation that many of these inequalities are redundant. For any $a_{1:i}$ with a proper (i.e. not full) prefix in $\min(\pa(S))$, we can add $a_{1:i}$ to $S$ without changing $\min(\pa(S))$, keeping the right hand side constant. This also increases the left hand side (since $R_\Phi$ is nonnegative). Thus, the strictest of these inequalities are precisely those where $S$ includes all paths with a proper prefix in $\min(\pa(S))$. For a fixed $T=\min(\pa(S))$, this makes $S$ the set of all $a_{1:j}$ where $a_{1:i} \in T$ for some $i<j$. Since all antichains are achievable as some $\min(\pa(S))$, this turns the above collection of inequalities into 
\begin{equation}
    \sum_{a_{1:i} \in T} \sum_{j=i+1}^{L+1} \sum_{a_{i+1:j} \in \V^{\leq j-i}} R_\Phi(a_{1:j}) \leq 1 - \left(1 - \sum_{a_{1:i} \in T} q(a_{1:i})\right)^K \quad \forall T \in \mathcal{A}(\V^{\leq L}).
\end{equation}
This covers the prefix-matching. The target-matching requirement remains the same as in \cref{lem:target}. Thus, all that remains is to show the objective values is the block efficiency. This holds as $R_\Phi(a_{1:i})$ is the probability that $\Phi$ outputs $a_{1:i}$, and the total number of tokens generated with this output is $i$, so summing over all possible outputs $a_{1:i} \in \V^{\leq L+1} \setminus \V^0$ gives the desired result.
\end{proof}

\end{document}
